\section*{The Appointed Texts: Rights-Limited Study Sheet}
\begin{studybox}{Text boundary}
The approved U.S. Lectionary and ICEL Missal English are protected and are not
reproduced. The Douay--Rheims (Challoner) is used only as the public-domain
study translation and is not the text proclaimed at Mass.
\end{studybox}
\subsection*{Ecclesiastes 1:2; 2:21--23}
The appointed cut is deliberately discontinuous. Its opening superlative
announces \emph{hevel}: something breath-like, transient, and resistant to
being mastered. It then resumes inside Qoheleth's experiment with toil. A
worker may combine wisdom, knowledge, skill, and sleepless care, yet cannot
control the character of the successor who receives the fruit. The point is
not that work has no proximate good, but that achievement cannot become a
permanent possession or answer the problem of death.
\subsection*{Psalm 90}
Human life returns to dust and passes like grass. The image moves within a
prayer addressed to the God who outlasts the generations: mortality becomes a
request for a wise heart, returning mercy, morning gladness, and stability for
the work of human hands. The last petition therefore prevents the first
reading from being heard as a command to abandon work.
The official occurrence page prints a response whose words derive from
Ps.~95:8 while marking it ``(1)'' beneath Psalm 90. This guide records the
source discrepancy and does not silently identify the response with Ps.~90:1.
\subsection*{Colossians 3:1--5, 9--11}
Those raised with Christ seek what is above because their life is hidden with
Christ and will be manifested with him. The cut then names disordered desires,
including greed as idolatry, before omitting vv.~6--8 and resuming with the
putting away of falsehood. The old person is stripped off and the new is being
renewed according to the creator's image; ethnic, ritual, cultural, and social
distinctions cannot become rival grounds of belonging where Christ is all and
in all.
\subsection*{Matthew 5:3 and Luke 12:13--21}
The acclamation blesses poverty of spirit and promises the Kingdom. In the
Gospel, a request that Jesus arbitrate an inheritance dispute is neither
decided nor endorsed. Jesus turns from the petitioner to the crowd, warns
against every form of greed, and distinguishes life from abundance. The
parable then exposes a prosperous man's inward monologue: land, produce,
barns, goods, and years all become possessives, but no neighbor or giver enters
his reckoning. God's address interrupts that closed account with death and the
question of who will receive the stores. The final contrast is between
treasuring for oneself and being rich toward God, not between negligence and
planning.
\subsection*{Missal formulary}
The Week XVIII owner controls the Entrance, three orations, and two
Communion-antiphon alternatives. Their protected English is not reconstructed
here; the Latin incipits and scriptural loci on page one identify them.
